\subsection{The No-U-Turn Sampler}
\label{subsec:nuts}

One of the hardest parts of tuning HMC is choosing the trajectory length.
The No-U-Turn Sampler (NUTS) of \citet[\S3.1]{hoffman_gelman_2014_nuts} automates this choice.
Its basic idea is simple: keep extending the HMC trajectory while it is still making productive forward progress, and stop once further integration would mostly retrace the same region.
The step size $\epsilon$ can still be tuned from acceptance behavior, but NUTS replaces the hand-chosen leapfrog count $L$ with a geometric stopping rule.
In a nearly elliptical geometry that rule tries to run long enough to cross the orbit without circling back and undoing earlier movement.

\begin{figure}[H]
  \centering
  \begin{tikzpicture}[>=Stealth]
  \path[use as bounding box] (-2.9,-2.55) rectangle (10.6,3.25);

  \begin{scope}[shift={(0,0)}]
    \node[font=\footnotesize] at (0,3.05) {still moving away};
    \draw[stat221Gray!50,line width=0.5pt] (0,0) ellipse (2.2 and 1.45);
    \draw[stat221Gray!35,dashed,line width=0.4pt] (0,0) ellipse (1.4 and 0.9);
    \coordinate (s1) at (-0.30,-1.43);
    \coordinate (c1) at (1.45,0.72);
    \fill[stat221Blue] (s1) circle (1.2pt);
    \fill[stat221Blue] (c1) circle (1.2pt);
    \draw[stat221Green!70!black,line width=1.0pt,->]
      (s1)
      .. controls (0.30,-1.25) and (1.10,-0.15) .. (c1);
    \draw[->,stat221Blue!80!black,line width=0.8pt] (s1) -- (c1);
    \draw[->,stat221Red!75!black,line width=0.9pt] (c1) -- ++(0.75,0.30);
    \node[stat221Gray!70,font=\footnotesize,anchor=east] at ($(s1)+(-0.06,-0.02)$) {$x_0$};
    \node[stat221Gray!70,font=\footnotesize,anchor=west] at ($(c1)+(0.08,0.02)$) {$x_t$};
    \node[stat221Blue!80!black,font=\footnotesize,anchor=west] at (0.30,-0.20) {$\Delta x$};
    \node[stat221Red!75!black,font=\footnotesize,anchor=west] at (2.15,1.05) {$M^{-1}p_t$};
    \node[font=\footnotesize] at (0,-2.15) {$\ip{\Delta x}{M^{-1}p_t} > 0$};
  \end{scope}

  \begin{scope}[shift={(7.8,0)}]
    \node[font=\footnotesize] at (0,3.05) {beginning to turn back};
    \draw[stat221Gray!50,line width=0.5pt] (0,0) ellipse (2.2 and 1.45);
    \draw[stat221Gray!35,dashed,line width=0.4pt] (0,0) ellipse (1.4 and 0.9);
    \coordinate (s2) at (-0.30,-1.43);
    \coordinate (c2) at (1.70,1.00);
    \fill[stat221Blue] (s2) circle (1.2pt);
    \fill[stat221Blue] (c2) circle (1.2pt);
    \draw[stat221Green!70!black,line width=1.0pt,->]
      (s2)
      .. controls (0.35,-1.20) and (1.85,-0.25) .. (2.05,0.65)
      .. controls (2.15,1.20) and (1.95,1.55) .. (c2);
    \draw[->,stat221Blue!80!black,line width=0.8pt] (s2) -- (c2);
    \draw[->,stat221Red!75!black,line width=0.9pt] (c2) -- ++(-0.85,0.10);
    \node[stat221Gray!70,font=\footnotesize,anchor=east] at ($(s2)+(-0.06,-0.02)$) {$x_0$};
    \node[stat221Gray!70,font=\footnotesize,anchor=west] at ($(c2)+(0.08,0.02)$) {$x_t$};
    \node[stat221Blue!80!black,font=\footnotesize,anchor=west] at (0.55,-0.10) {$\Delta x$};
    \node[stat221Red!75!black,font=\footnotesize,anchor=east] at (0.95,1.35) {$M^{-1}p_t$};
    \node[font=\footnotesize] at (0,-2.15) {$\ip{\Delta x}{M^{-1}p_t} < 0$};
  \end{scope}
\end{tikzpicture}

  \caption{A schematic of the NUTS stopping idea.
  While the secant vector $\Delta x$ and the velocity $M^{-1}p$ point in roughly the same direction, extending the trajectory is still productive.
  Once their inner product becomes negative, the trajectory is turning back toward where it came from and further leapfrog steps are likely to waste computation.}
  \label{fig:nuts-uturn}
\end{figure}

\subsubsection{The U-turn stopping criterion}
\label{subsubsec:nuts-uturn}

Suppose an exact Hamiltonian trajectory starts at $x_0$ and reaches $x_t$ with momentum $p_t$.
If we write $\Delta x_t \coloneqq x_t-x_0$, then the squared distance from the starting point changes according to
\begin{equation}
  \frac{d}{dt}\,\frac{1}{2}\norm{x_t-x_0}^2
  =
  \ip{x_t-x_0}{\dot x_t}
  =
  \ip{x_t-x_0}{M^{-1}p_t}.
  \label{eq:nuts-distance-derivative}
\end{equation}
Equation \eqref{eq:nuts-distance-derivative} is the key geometric diagnostic.
As long as it is positive, the trajectory is still moving away from its starting point.
When the derivative becomes negative, the trajectory has begun to turn back.
On an elliptical orbit this is exactly the moment at which continuing the simulation starts to undo previous movement rather than discovering new regions of the typical set.

This observation gives the core stopping idea: once the trajectory begins to turn back, additional leapfrog steps are no longer buying new exploration.
The probabilistic difficulty is that we cannot simply stop at the first realized U-turn and keep that endpoint, because the stopping rule itself depends on the realized path.
The remaining construction packages this geometric stopping idea into a symmetric HMC transition.

\subsubsection{Symmetric tree expansion and candidate selection}
\label{subsubsec:nuts-tree}

NUTS implements that stopping idea with two pieces of machinery.
First, a slice variable excludes states with excessively large energy.
Second, a doubling construction grows a symmetric set of candidate states around the current point instead of privileging a single realized endpoint.

The slice variable
\begin{equation}
  u \sim \mathrm{Unif}\!\left(0,e^{-H(x_0,p_0)}\right),
  \label{eq:nuts-slice}
\end{equation}
where $(x_0,p_0)$ is the current phase-space state after momentum refreshment.
Only states satisfying
\[
  e^{-H(x,p)} \ge u
\]
are eligible candidates for the next draw.
This slice restriction prevents obviously poor numerical states from entering the candidate set and makes the correctness argument symmetric.

Starting from the current state, NUTS builds a balanced binary tree by repeatedly choosing a direction $v_j\in\{-1,+1\}$ and integrating $2^j$ leapfrog steps of size $v_j\epsilon$ at depth $j$.
At every stage it maintains left and right endpoints $(x^-,p^-)$ and $(x^+,p^+)$, along with the candidate states encountered so far that satisfy the slice condition.
Concretely, depth \(0\) adds one leapfrog step to one side of the current state.
If no stop is triggered, depth \(1\) appends a fresh subtree of length \(2\), then depth \(2\) appends one of length \(4\), and so on, so the explored path length doubles at each successful stage.
The tree is therefore just a bookkeeping device for this recursive concatenation.
One can think of each subtree as carrying a small summary: its left and right endpoints, the number of slice-admissible states inside it, and a flag saying whether a U-turn or divergence has already appeared.
Those summaries are enough to keep the construction symmetric between forward and backward time without privileging the order in which individual leapfrog states were generated.
This bookkeeping also explains candidate selection.
After a valid subtree has been grown, NUTS does not simply keep the newest endpoint.
Instead it merges the old and new subtrees and updates the candidate by sampling from the union with probability proportional to the number of slice-admissible states contributed by each side.
That recursive size-weighted merge is how implementations reproduce the uniform law on the admissible candidate set without explicitly storing every visited point in a long trajectory.

Tree expansion stops as soon as one of two things happens.
The first is a divergence, meaning the energy error has become too large for the numerical trajectory to be trusted.
The second is an endpoint version of the U-turn condition:
\begin{equation}
  \ip{x^+ - x^-}{M^{-1}p^-} < 0
  \qquad\text{or}\qquad
  \ip{x^+ - x^-}{M^{-1}p^+} < 0.
  \label{eq:nuts-endpoint-criterion}
\end{equation}
The secant $x^+-x^-$ summarizes the currently explored trajectory segment, and the two endpoint momenta check whether either end is now pointing back across that segment.
When \eqref{eq:nuts-endpoint-criterion} holds, further doubling would mostly revisit the same region rather than extend exploration.
\Cref{fig:nuts-tree-growth} illustrates the doubling logic schematically.
There is also no separate Metropolis accept/reject step at the end of this construction.
Energy fluctuations are handled indirectly by the slice restriction and by terminating the expansion when a divergence is detected, after which the next state is drawn from the admissible set already built inside the slice.

\begin{figure}[H]
  \centering
  \begin{tikzpicture}[>=Stealth]
  \path[use as bounding box] (-4.1,-2.55) rectangle (5.05,2.85);

  \begin{scope}
    \draw[stat221Gray!40,line width=0.45pt] (0,0) circle (1.10);
    \draw[stat221Gray!40,line width=0.45pt] (0,0) circle (1.65);
    \draw[stat221Gray!40,dashed,line width=0.45pt] (0,0) circle (2.15);
    \path[fill=stat221Blue!12, even odd rule] (0,0) circle (2.15) (0,0) circle (1.45);

    \coordinate (r0)  at (-2.05,-0.15);
    \coordinate (r1a) at (-1.55,0.85);
    \coordinate (r1b) at (-1.52,-1.00);
    \coordinate (r2a) at (-0.95,1.45);
    \coordinate (r2b) at (-0.70,0.35);
    \coordinate (r2c) at (-0.88,-0.42);
    \coordinate (r2d) at (-0.45,-1.55);

    \coordinate (p0) at (202:1.65);
    \coordinate (p1a) at (170:1.65);
    \coordinate (p1b) at (118:1.65);
    \coordinate (p2a) at (152:1.65);
    \coordinate (p2b) at (92:1.65);
    \coordinate (p2c) at (58:1.65);
    \coordinate (p3a) at (138:1.65);
    \coordinate (p3b) at (106:1.65);
    \coordinate (p3c) at (72:1.65);
    \coordinate (p3d) at (38:1.65);

    \draw[stat221Gray!75,line width=0.9pt] (r0) -- (r1a);
    \draw[stat221Gray!75,line width=0.9pt] (r0) -- (r1b);
    \draw[stat221Gray!75,line width=0.9pt] (r1a) -- (r2a);
    \draw[stat221Gray!75,line width=0.9pt] (r1a) -- (r2b);
    \draw[stat221Gray!75,line width=0.9pt] (r1b) -- (r2c);
    \draw[stat221Gray!75,line width=0.9pt] (r1b) -- (r2d);

    \draw[stat221Gray!55,line width=0.85pt] (r2a) -- (p3a);
    \draw[stat221Gray!55,line width=0.85pt] (r2a) -- (p3b);
    \draw[stat221Gray!55,line width=0.85pt] (r2b) -- (p2b);
    \draw[stat221Gray!55,line width=0.85pt] (r2b) -- (p2a);
    \draw[stat221Gray!55,line width=0.85pt] (r2c) -- (p1a);
    \draw[stat221Gray!55,line width=0.85pt] (r2c) -- (p1b);
    \draw[stat221Gray!55,line width=0.85pt] (r2d) -- (p3c);
    \draw[stat221Gray!55,line width=0.85pt] (r2d) -- (p3d);

    \draw[stat221Gray!45,dashed,line width=0.55pt] (210:1.65) arc[start angle=210,end angle=24,radius=1.65];

    \fill[black] (p0) circle (1.8pt);
    \fill[stat221Blue] (p1a) circle (2.0pt);
    \fill[stat221Blue] (p1b) circle (2.0pt);
    \fill[stat221Green!75!black] (p2a) circle (2.0pt);
    \fill[stat221Green!75!black] (p2b) circle (2.0pt);
    \fill[stat221Green!75!black] (p2c) circle (2.0pt);
    \fill[stat221Purple] (p3a) circle (2.0pt);
    \fill[stat221Purple] (p3b) circle (2.0pt);
    \fill[stat221Purple] (p3c) circle (2.0pt);
    \fill[stat221Purple] (p3d) circle (2.0pt);

    \node[stat221Gray,font=\footnotesize,anchor=west] at (1.95,1.42) {level curves of $U(x)$};
    \node[stat221Gray,font=\footnotesize,anchor=west] at (1.95,1.04) {doubling expands the tree};

    \matrix[
      anchor=west,
      draw=stat221Gray!25,
      fill=white,
      rounded corners,
      inner sep=2.5pt,
      row sep=0.10cm,
      column sep=0.16cm,
      font=\footnotesize,
    ] at (-3.85,1.72) {
      \node{\tikz[baseline=-0.5ex]\fill[black] (0,0) circle (1.2pt);}; & \node{$j=0$}; \\
      \node{\tikz[baseline=-0.5ex]\fill[stat221Blue] (0,0) circle (1.7pt);}; & \node{$j=1$}; \\
      \node{\tikz[baseline=-0.5ex]\fill[stat221Green!75!black] (0,0) circle (1.7pt);}; & \node{$j=2$}; \\
      \node{\tikz[baseline=-0.5ex]\fill[stat221Purple] (0,0) circle (1.7pt);}; & \node{$j=3$}; \\
    };
  \end{scope}
\end{tikzpicture}

  \caption{A schematic NUTS tree grown along an HMC trajectory.
  Each doubling adds a new half-tree of leapfrog states, increasing the explored path length exponentially until a U-turn or a divergence stops the expansion.}
  \label{fig:nuts-tree-growth}
\end{figure}

\begin{algorithm}[H]
  \caption{High-level NUTS iteration}
  \label{alg:nuts}
  \begin{algorithmic}[1]
    \Require Current state $x^{(t)}$, mass matrix $M$, step size $\epsilon$, maximum tree depth $J_{\max}$
    \State Draw $p_0 \sim \mathcal{N}(0,M)$ and $u \sim \mathrm{Unif}\!\left(0,e^{-H(x^{(t)},p_0)}\right)$.
    \State Initialize left and right endpoints at $(x^{(t)},p_0)$ and candidate set $\mathcal{C}\gets\{(x^{(t)},p_0)\}$.
    \For{$j=0,1,\dots,J_{\max}$}
      \State Choose a direction $v_j\in\{-1,+1\}$.
      \State Extend the corresponding endpoint by $2^j$ leapfrog steps of size $v_j\epsilon$.
      \State Add any new states satisfying the slice condition $e^{-H(x,p)}\ge u$ to $\mathcal{C}$.
      \If{a divergence occurs or the endpoint criterion \eqref{eq:nuts-endpoint-criterion} holds}
        \State \textbf{break}
      \EndIf
    \EndFor
    \State Sample the next state $x^{(t+1)}$ using the candidate-selection rule that preserves the uniform law on admissible candidates in $\mathcal{C}$.
\end{algorithmic}
\end{algorithm}

The original algorithm samples uniformly from the admissible candidate set inside the slice.
Modern implementations often use a closely related multinomial sampling rule that weights states by $e^{-H(x,p)}$ instead of enforcing the slice variable explicitly.
The implementation details differ, but their role is the same: turn the geometric rule ``stop when the trajectory starts to retrace'' into a valid symmetric transition kernel.
The maximum tree depth \(J_{\max}\) is only a computational cap on this doubling procedure.
Hitting that cap means the sampler would have kept exploring if allowed; it does not indicate a correctness failure, only that the implementation has limited the longest admissible trajectory for cost reasons.

\subsubsection{Why the transition is still valid}
\label{subsubsec:nuts-correctness}

The key point is that NUTS is \emph{not} ``run HMC until the first U-turn and keep the endpoint.''
That naive rule would privilege one realized stopping point and would generally break the symmetry needed for detailed balance.
Instead, NUTS uses the slice variable and the doubling tree to construct a candidate set that contains the current state, excludes obviously bad high-energy states, and is built in a way that is symmetric between forward and backward integration; this is the core construction in \citet[\S\S3.1.1--3.1.2]{hoffman_gelman_2014_nuts}.

Once that symmetric candidate set has been built, the next state is drawn by a transition rule that leaves the uniform distribution on the admissible candidates invariant.
In the simplified presentation that means sampling uniformly from the slice-restricted set.
In the efficient implementation it means using subtree sizes to reproduce the same uniform law without storing every visited state explicitly.
The correctness mechanism is therefore easy to summarize even though the implementation details are intricate:
NUTS preserves HMC's reversible geometry by replacing a single path-dependent endpoint with a symmetric set of admissible states and then sampling from that set with a reversible rule.

\subsubsection{Warm-up and practical tuning}
\label{subsubsec:nuts-warmup}

In practice NUTS is almost always paired with automatic adaptation of the step size and mass matrix during warm-up.
Warm-up has two jobs: choose a step size that gives stable, high-acceptance trajectories, and estimate a mass matrix that matches the scale of the target.
The original NUTS paper uses dual averaging on $\log\epsilon$ to target a chosen average acceptance statistic $\delta$; see \citet[\S3.2, p.~1608]{hoffman_gelman_2014_nuts}.
In the optimization language of \cref{subsubsec:dual-averaging}, the quantity being updated is now the scalar tuning parameter $\log\epsilon$ rather than the model parameter itself, and the role of the noisy gradient is played by the acceptance error.
If $\alpha_j$ denotes the acceptance statistic from warm-up iteration $j$, one defines $H_j\coloneqq \delta-\alpha_j$ and updates
\[
  \bar H_j
  =
  \left(1-\frac{1}{j+t_0}\right)\bar H_{j-1}
  +
  \frac{1}{j+t_0}H_j,
  \qquad
  \log \epsilon_j
  =
  \mu - \frac{\sqrt{j}}{\gamma}\bar H_j,
\]
for a baseline value $\mu$, then forms the smoothed average
\[
  \log \bar\epsilon_j
  =
  j^{-\kappa}\log\epsilon_j
  +
  \bigl(1-j^{-\kappa}\bigr)\log \bar\epsilon_{j-1}.
\]
The first recursion aggregates acceptance errors, while the second maps that accumulated dual signal back to the primal tuning variable $\log\epsilon$.
The step-size recursion is therefore also called \emph{Nesterov's primal-dual averaging}: it maintains an aggregated dual error signal and maps it back to a primal tuning variable.
It is therefore distinct from Nesterov's accelerated gradient, which updates the model parameter by a look-ahead gradient step.
The commonly used defaults from that construction are $\gamma=0.05$, $t_0=10$, and $\kappa=0.75$.

Warm-up is staged because the two tuning problems live on different time scales.
The step size reacts quickly: after a short run it is already clear whether trajectories are numerically stable and whether the acceptance statistic is near the target.
The mass matrix is slower because it requires a covariance estimate from a window of warm-up draws that is long enough to say something reliable about scale.

A typical workflow therefore has three phases.
One starts with \(M=I\) (or a simple diagonal guess) and tunes only \(\epsilon\) until trajectories are numerically stable and the average acceptance statistic is in a sensible range.
Then come the slower covariance windows: warm-up draws are accumulated, a new estimate of \(M^{-1}\) is formed, and the step-size adaptation is restarted because the geometry has just changed.
This middle stage is repeated with longer windows until the metric estimate stops moving enough to matter.
Only then does one freeze \(M\) and do a final short retuning pass for \(\epsilon\) with the geometry held fixed.
This is the practical meaning of the staged warm-up used in modern NUTS implementations: fast adaptation for stability, slower windowed adaptation for geometry, then fast retuning once the geometry has been fixed.
The common summary is therefore ``fast-slow-fast'': short early adaptation for \(\epsilon\), long middle windows for \(M\), then a short final cleanup pass for \(\epsilon\) with \(M\) fixed.

As in HMC, the geometric rationale for the metric comes from \Cref{subsec:geometry-control}; the sampler-specific point is that a better metric makes the no-U-turn rule act on a trajectory whose numerical scale is already better balanced across directions.
Because the covariance estimates are noisy, practical implementations use windowed updates, diagonal approximations, or mild shrinkage and ridge regularization before freezing the metric.
Once warm-up is finished, the tuning parameters are frozen and standard post-warm-up diagnostics such as $\hat R$, effective sample size, trace plots, and divergence counts become the relevant checks; see \Cref{subsec:mcmc-diagnostics}.

From a practical point of view, NUTS removes the most fragile hand-tuned HMC parameter.
After warm-up, the user no longer chooses $L$ directly: the realized trajectory length is whatever the no-U-turn logic allows before retracing begins.
Modern probabilistic programming systems therefore often expose NUTS, rather than plain fixed-$L$ HMC, as the default general-purpose gradient-based sampler.
